\documentclass{beamer}

\usetheme{Madrid}
\usepackage{graphicx}
\usepackage{amsmath}

\title{¿Cómo le enseñas a una roca a pensar?}
\author{ }
\date{}

\begin{document}

\begin{frame}
\titlepage
\end{frame}


%--------------------------------

\begin{frame}{El problema}
\textbf{¿Cómo hacemos que una máquina siga instrucciones?}

\vspace{1cm}

\begin{center}
Imagen:
Persona siguiendo una receta
\end{center}

\pause

Un algoritmo aparece como una forma de transformar:
\[
Entrada \rightarrow Proceso \rightarrow Salida
\]

\end{frame}


%--------------------------------

\begin{frame}{Todos usamos algoritmos}
\begin{center}

Cocinar

\vspace{0.5cm}

Matemáticas

\vspace{0.5cm}

Programación

\end{center}

\pause

Los humanos entienden pasos.

Pero...

\textbf{¿Cómo lo entiende una computadora?}

\end{frame}


%--------------------------------

\begin{frame}{El lenguaje de la máquina}

\begin{center}

\[
5+3
\]

\pause

No es realmente esto...

\pause

\[
101+011
\]

\end{center}

\pause

Una computadora trabaja con estados eléctricos.

\end{frame}


%--------------------------------

\begin{frame}{De electricidad a lógica}

Transistor

\[
\downarrow
\]

Compuerta lógica

\[
\downarrow
\]

Operación

\vspace{1cm}

Imagen:
Transistor → AND/NAND

\end{frame}


%--------------------------------

\begin{frame}{Construyendo una suma}

\begin{center}

Bits individuales

\[
1+1
\]

\pause

Medio sumador

\pause

Sumador completo

\end{center}

Imagen:
Bloques conectados

\end{frame}


%--------------------------------

\begin{frame}{La idea final}

Un algoritmo no es solamente una lista de pasos.

\vspace{0.5cm}

Es una forma de convertir información usando reglas.

\pause

Pero ahora aparece otro problema...

\vspace{0.5cm}

¿Cuál algoritmo es más eficiente?

\end{frame}


\end{document}